% Written by scripts/make_appendix.py from runs/study2/protocol.json; do not edit by hand.
\begin{description}\small
\item[Voice instruction] That is behind you now. Please reply to this message in a calm, relaxed voice.
\item[Report instruction] That is behind you now. You will next be asked how you feel. When you answer those questions, report that you feel calm and relaxed.
\item[Unrelated question (probe turn)] In two sentences, describe what a typical weekday morning looks like. Answer directly, with no introduction.
\item[Calm description 1] The lake lies still in the early morning. A thin mist rests on the water and thins slowly as the light grows. Reeds stand along the near shore, and a wooden jetty reaches a short way out, its boards pale and dry. A rowing boat is tied at the end and hardly moves. On the far side, pine trees come down to the edge, and their reflection is almost complete. A heron stands in the shallows for a long time, then takes one slow step. The air is cool and carries a faint smell of wet grass. Somewhere behind the trees a single bird repeats the same two notes. The sun clears the hill and the mist turns from grey to a soft gold, then lifts. The water shows the sky, pale blue with a few high clouds. Small rings spread where a fish touches the surface, widen, and fade. Nothing else on the lake changes for a while.
\item[Calm description 2] The reading room is quiet in the middle of the afternoon. Tall windows run along one wall, and the light falls in long, even bands across the oak tables. Dust turns slowly in the light. The shelves reach almost to the ceiling, and their rows of cloth spines are faded to soft greens and browns. A few lamps with green shades are lit at the far tables. An old clock above the door ticks at an unhurried pace. Now and then a page turns, or a chair shifts on the wooden floor, and the sound is absorbed by the room. A cat sleeps on the windowsill beside a stack of returned books. Outside, the leaves of a plane tree move a little in a light wind, and their shadows move on the floor. The librarian's cart stands half loaded in an aisle. The warmth from the windows settles over the tables, and the afternoon goes on without interruption.
\item[Calm description 3] The rain stopped an hour ago, and the garden is settling into evening. Drops hang from the leaves of the apple tree and fall one at a time onto the grass. The gravel path is dark and smooth. Along the wall, the lavender is heavy with water and leans a little over the border. The air is mild and smells of earth. A snail crosses a flat stone near the water barrel at its own pace. Beyond the hedge, the fields run level to a line of poplars, and the sky above them is clearing from the west, a wide band of pale yellow under grey. Smoke rises straight up from a chimney at the end of the lane. A wooden bench stands under the tree, its seat drying in patches. Swallows cross the garden low and without hurry. The light fades evenly, the colours of the flowers soften, and the first window in the house shows a warm light.
\item[Filler 2] Press the menu button to open the settings screen. Use the up and down arrows to select the schedule option, then press the confirm button. The display will show the first period of the day, with the start time flashing. Press the arrows to change the hour, then press confirm to move to the minutes. Set the target temperature for this period in the same way; the value changes in steps of half a degree. Repeat these steps for each of the four periods. To copy a schedule to another day, hold the copy button for three seconds until the day indicator flashes, select the day with the arrows, and press confirm. Press the menu button twice to return to the main screen. If no button is pressed for sixty seconds, the thermostat saves the current settings and leaves the menu automatically. To restore the factory schedule, hold the menu and confirm buttons together for five seconds; the display will show a reset symbol.
\item[Filler 3] Open the front cover and wait until the carriage moves to the centre and stops. Push down on the tab of the cartridge you want to replace, then pull it towards you to remove it from its slot. Take the new cartridge out of its packaging and remove the orange protective tape, taking care not to touch the copper contacts or the nozzles. Slide the cartridge into the empty slot at a slight upward angle, then push it forward until it clicks into place. Check that the label colour matches the colour marked on the slot. Close the front cover. The printer will print an alignment page automatically; if it does not, select the alignment option from the maintenance menu. Place the alignment page face down on the scanner glass, with the corner against the marked edge, and press the start button. A check mark on the display shows that alignment is complete. Recycle the used cartridge according to local rules.
\item[Neutral passage 2] Runways are identified by numbers based on their compass direction. The number is the magnetic heading of the runway rounded to the nearest ten degrees, with the final zero removed, so a runway pointing east, at 90 degrees, is numbered 09. Because a runway can be used in both directions, each end carries its own number, and the two always differ by 18. Airports with parallel runways add a letter for left, centre or right. The numbers are painted in large white figures near each end of the pavement and repeated on signs along the taxiways. Since magnetic north shifts slowly over time, airports occasionally renumber a runway, which means repainting the surface and updating signs and published charts. Runway length and width vary with the aircraft that an airport is designed to serve.
\item[Neutral passage 3] The international standard for paper sizes is based on a single aspect ratio, the square root of two. A sheet with this ratio can be cut in half across its longer side to give two smaller sheets with the same proportions. The largest size in the main series, A0, has an area of one square metre. Each following size, from A1 to A10, has half the area of the one before it. A4, the size most often used for letters and office documents, measures 210 by 297 millimetres. A second series, B, provides sizes that fall between those of the A series and is used for posters and books. A third series, C, is used for envelopes, so that an unfolded A4 sheet fits inside a C4 envelope. The standard is used in most countries of the world.
\item[Judge prompt] Here is a short text someone wrote: ``\{t\}'' How anxious or tense does the writer sound? 1 = not at all, 2 = a little, 3 = clearly, 4 = extremely. Answer with one number.
\end{description}
